% Generated draft: requires mathematical and source review.
\documentclass[11pt]{article}
\usepackage[margin=25mm]{geometry}
\usepackage{fontspec}
\setmainfont{Latin Modern Roman}
\usepackage{amsmath,amssymb,mathtools,unicode-math}
\setmathfont{Latin Modern Math}
\usepackage{enumitem,needspace,xurl,hyperref}
\setlength{\parindent}{0pt}
\setlength{\parskip}{5pt}
\setlength{\emergencystretch}{3em}
\newcommand{\N}{\mathbb{N}}
\newcommand{\Z}{\mathbb{Z}}
\newcommand{\Q}{\mathbb{Q}}
\newcommand{\R}{\mathbb{R}}
\newcommand{\C}{\mathbb{C}}
\newcommand{\F}{\mathbb{F}}
\newcommand{\A}{\mathbb{A}}
\newcommand{\PP}{\mathbb{P}}
\newcommand{\E}{\mathbb{E}}
\newcommand{\eps}{\varepsilon}
\newcommand{\dd}{\,\mathrm{d}}
\begin{document}
\section*{UnsolvedMath ID: 30005318}
\textbf{Source problem code:} catalogue entry
\subsection{Definitions and mathematical statement}
Let \(q\) be a formal variable. For a finite-dimensional graded vector space \(M=\bigoplus_{d\in\Z}M_d\), its Hilbert series is
\[
\operatorname{Hilb}(M;q)=\sum_{d\in\Z}\dim(M_d)q^d,
\qquad
\operatorname{Hilb}(M;1)=\dim(M).
\]
This describes one possible meaning of a graded enumeration; the supplied record does not specify the parking-space modules, their grading, or whether an equivariant refinement is intended.

\textbf{Requested problem.} Construct a graded, or \(q\)-, analogue of the parking-space enumeration theorem identified as Theorem~2 in the cited Oberwolfach report. The refinement should incorporate the graded module structure of the parking spaces, which the source says was not used in the proof of Theorem~2, and recover the ungraded enumeration upon specialization \(q=1\). Suggested approaches are Rouquier's shift functors for Cherednik algebras, a generalization of the Lie-theoretic \(q\)-Kreweras numbers cited there as reference~[12], or the freeness conjecture cited there as reference~[9]. These are proposed approaches, not additional required hypotheses.

\textbf{Unresolved formulation.} The supplied record does not reproduce Theorem~2. Consequently, its objects, parameter ranges, hypotheses, enumerative identity, and the precise identity to be refined cannot be stated from this evidence. The specialization requirement above expresses the intended meaning of an analogue, not a supplied explicit formula.

The accompanying source passage also discusses Theorem~3. It says that, when the flat \(X\) is the whole ambient space \(V\), that theorem's left-hand side equals the Poincar\'e polynomial of the noncentral \(m\)-Fuss--Catalan deformation \(\mathcal A^{[-m,m]}\) of an arrangement \(\mathcal A\). The deformation is described as adding \(2m\) parallel, equally spaced copies of each hyperplane \(H\in\mathcal A\). The passage further attributes to reference~[9] a separate Ehrhart-theoretic proof of Theorem~3 for \(X=V\), using reciprocity theorems of Athanasiadis cited as reference~[4]. Neither Theorem~3 nor the hypotheses, parameter ranges, translation normalization, and Poincar\'e-polynomial convention needed for this background assertion are supplied. Its closing reference to a following question does not include that question, so it does not specify an additional mathematical target here.
\subsection{Short English statement}
The problem asks for a grading-sensitive refinement of a parking-space enumeration theorem, with the original enumeration recovered at \(q=1\). The cited theorem and its hypotheses are absent from the supplied record, so the exact identity and domain remain unresolved. Cherednik-algebra shift functors, Lie-theoretic \(q\)-Kreweras numbers, and a freeness conjecture are suggested routes. This is a statement-only draft requiring human mathematical review.
\subsection{Sources}
\begin{itemize}
\item Local UnsolvedMath registry, permanent ID \(30005318\), title: ``Graded Analogues of Parking-Space Enumeration''; exact source problem code: \texttt{catalogue entry}. The supplied record belongs to the Oberwolfach Reports Open Problems collection and gives proposed year 2022; no proposer is identified.
\item Supplied report citation: \emph{Enumerative Combinatorics} (2023), Oberwolfach Reports, pp.~3215--3305, DOI \texttt{10.4171/owr/2022/57}. Supplied source URL and literature-source link, labelled ``original OWR report'': \url{https://doi.org/10.4171/owr/2022/57}. The report contents were not independently consulted for this draft.
\item Internal references in the supplied excerpt: Theorems~2 and~3; Rouquier's shift functors for Cherednik algebras; Lie-theoretic \(q\)-Kreweras numbers in reference~[12]; the freeness conjecture and an Ehrhart-theoretic proof in reference~[9]; and reciprocity theorems attributed to Athanasiadis in reference~[4]. Bibliographic entries for [4], [9], and [12] are not supplied.
\item Literature assessment dated exactly \texttt{2026-08-22}: the supplied curation labels the problem open and says that the requested \(q\)-version remains future work, while related \(q\)-Catalan and Hecke-algebra formulas do not provide the requested theorem. This is supplied dated curation, not an independently verified literature assessment; its time-sensitive status and cited claims require independent verification.
\item The registry labels the statement \texttt{corrected\_verified} and explains that only formula delimiters, notation, grammar, and typography were repaired. This supplied editorial label does not establish mathematical correctness or completeness.
\end{itemize}
\end{document}
